\newpage
\newcommand{\ResearchScaleError}{4.4e-16}
\newcommand{\ResearchPass}{70}
\newcommand{\ResearchFollowPass}{8}
\newcommand{\ResearchFollowChange}{0.0011\%}
\newcommand{\ResearchCorrection}{2.288\%}
\newcommand{\ResearchQuadError}{4.4e-12}
\newcommand{\ResearchQuadTail}{7.5e-12}

\eyebrow{RESEARCH FOLLOW-UP / SEPTEMBER 10, 2026}
\section*{1. Does the solver preserve cubic scaling?}
The literature review suggested replacing a purely visual search with experiments
that target the structure of the dispersion. We begin with an exact symmetry of
the homogeneous equation, then vary spectral geometry and cancellation, and test
an independent integral calculation.
\[
 f_\lambda(X)=\lambda^3 f(\lambda X),\qquad \|f_\lambda\|_1=\|f\|_1,
 \qquad R(t;f_\lambda)=R(\lambda^3t;f).
\]
Here $X=(x,y,z)$ and $\beta=0$. At $t=8/\lambda^3$, all three rescaled
inputs should produce the same compensated ratio as the original at $t=8$.
\pic{1}{visuals/research_scaling.pdf}
\takeaway{All eleven recipes pass the scaling check. Across 33 evaluations,
the largest relative complex-field discrepancy is \texttt{\ResearchScaleError}.}
The calculation uses $\lambda=1/2,1,2$, hence times $64,8,1$. A length-144,
$96^3$ grid is rescaled to length $144/\lambda$ and spacing $1.5/\lambda$.
Initial data and the dispersion symbol are reevaluated on each scaled grid.
The left panel shows three representative recipes; the right summarizes all eleven.

\textbf{What this verifies.} This is a check of the discrete implementation of an
exact symmetry. The matched grids make the relation exact even for a finite
periodic model; machine-precision agreement is not a new whole-space error bound.

\textbf{What scaling says about the research question.} If a bound with exponent
$\gamma$ and the same $L^1$ right-hand side held uniformly for every input at all
$t\ge1$, then at $t=\lambda^{-3}$ it would require
\[
 \lambda^{7/2}\big\||D_x|^{1/2}S(1)f\big\|_\infty
 \le C\lambda^{3\gamma}\|f\|_1.
\]
For a nonzero weighted solution this is impossible as $\lambda\to0$ when
$\gamma>7/6$. Improving the homogeneous result therefore requires a different
question: a better constant, or faster decay under additional input assumptions.
This derivation does not apply unchanged to fixed $\beta=1$.

\newpage
\eyebrow{RESEARCH FOLLOW-UP / TARGETED FREQUENCIES}
\section*{2. Where the spectrum sits changes the ratio}
We vary longitudinal width $\sigma_x\in\{3,6\}$ and spectral offset
$q_x+\beta\in\{-0.75,-0.25,0,0.25,0.75\}$, holding transverse width at 3.
These give twenty cases. Three more target the second curvature-degenerate set,
two impose zero spatial integral, and one supplies a Gaussian cancellation control:
26 inputs, each evaluated at $t=4,8,16$.
\pic{.94}{visuals/research_sweep.pdf}
\textbf{Reading the markers.} Circles use the fine 3D FFT result after passing
the 2\% screen. Diamonds use a separate Gaussian-integral follow-up after the
initial screen failed. The dotted line marks the multiplier's cusp.

\takeaway{The homogeneous ratios are largest at the central frequency among
these sampled offsets. The inhomogeneous patterns differ. Frequency location
and width both matter; there is no single ordering that transfers across regimes.}
The original checks compare lengths 144 and 192 at spacing 1.5, then spacings
1.5 and 1 at length 192. \ResearchPass\ of 78 case--time measurements pass the
2\% screen. Eight are flagged and recomputed using analytic transverse integration,
periods 8192 and 16384, radial limits 256 and 512, and local integral optimization. For the odd Gaussian control, the Fourier
amplitude is multiplied by $-i\sqrt{\pi/2}\,\sigma_x k$ to enforce exact continuum
$L^1=1$ normalization; the other follow-ups use normalized Gaussians.
\ResearchFollowPass\ of eight pass the 0.1\% follow-up screen. The largest change
between those two broad searches is \ResearchFollowChange; the largest correction
to an original fine-grid peak is \ResearchCorrection.

\textbf{Second degeneracy.} For the inhomogeneous symbol we also test
$q_x=0.5$, $q_y=\sqrt{0.75}+\delta$ with $\delta=-0.2,0,0.2$ and widths 6.
The middle center lies on $q_x(q_x+1)=q_y^2$. These are localized Gaussian spectra,
not exact support on the surface. All three pass the original screen at $t=16$;
the central case is not the largest of the three. This small test does not identify
a worst-case input.

\newpage
\eyebrow{RESEARCH FOLLOW-UP / CANCELLATION AND CONCENTRATION}
\section*{3. Test assumptions, not just different shapes}
For each regime, compare an unshifted Gaussian with
$f(X)\propto x\exp[-|X|^2/18]$, normalized separately to $L^1=1$.
The latter has zero spatial integral by oddness: low-frequency cancellation is
built into the input rather than inferred from a picture.
\pic{.97}{visuals/research_cancellation.pdf}
\textbf{Interpretation.} These are short-time ratios, not fitted asymptotic rates.
The inhomogeneous zero-integral ratio is almost flat from $8$ to $16$, within
the numerical differences tested here. A zero integral changes the curve but does not, by itself, establish a stronger
uniform estimate. These inputs still have Gaussian spectral tails; no experiment
with strictly supported spectra away from degeneracies has been completed.
\pic{.97}{visuals/research_quadrature.pdf}
\textbf{Left: concentrated inputs.} At $t=1$, exactly $L^1$-normalized Gaussian
widths decrease from 3 to 0.75. Ratios increase in both regimes. All ten cases pass
a 2\% numerical screen: longitudinal period/radial extent are doubled, radial
spacing is halved, and candidate peaks are checked by non-periodic integration.
The longitudinal spacing is 0.25 and the narrowest width spans three such steps.

\takeaway{Narrower inputs probe larger ratios here, but five widths do not
establish convergence to a kernel maximum or determine the optimal constant.}
\textbf{Right: independent integration.} The comparison is explained on the next
page. Its later-time points evaluate fixed coordinates, not newly located maxima.

\newpage
\eyebrow{RESEARCH FOLLOW-UP / INDEPENDENT INTEGRATION}
\section*{4. A second quadrature agrees through $t=1024$}
We compare two separately implemented methods for the reduced Gaussian integral.
Both handle the multiplier cusp at $k=-\beta$, but use different panel placement
and integration rules.
\begin{center}\begin{tabular}{p{.26\linewidth}p{.65\linewidth}}\toprule
Method & How it evaluates the integral\\\midrule
Fixed Gauss--Legendre & 64 Gaussian-frequency panels plus a cusp split;
$k=-\beta\pm u^2$ regularization; 256 and 512 nodes per panel.\\
Adaptive Gauss--Kronrod & Panels limit the principal cubic phase variation;
stationary points of $tk^3+xk$ are inserted; the residual complex amplitude is
integrated adaptively after a separately implemented cusp substitution.\\\bottomrule
\end{tabular}\end{center}
\takeaway{Across twelve pointwise comparisons, the maximum relative complex-field
difference is \texttt{\ResearchQuadError}. No adaptive integration warnings were
recorded. Expanding frequency truncation from 12 to 14 Gaussian standard deviations
changes values by at most \texttt{\ResearchQuadTail} relatively.}
For each of the four earlier Gaussian-family cases, we keep its refined $t=256$
peak coordinates and evaluate the field there at $256,512,1024$. The 512-node
Gauss--Legendre values agree with the adaptive calculation. For two cases at
$1024$, changing from 256 to 512 nodes changes the field by about $10^{-7}$
relatively, showing why integration order needs checking as time increases.

\textbf{What has advanced.} We have verified a scaling identity, completed a
controlled frequency sweep with follow-ups for every initially flagged point,
compared cancellation inputs, tested a family of concentrated inputs, and checked
a second integration implementation at later times.

\textbf{What has not advanced yet.} Global peak searches at $512$ and $1024$
have not been performed. Their locations can move, so the pointwise comparisons
cannot extend the decay curves. The new screens remain numerical consistency
checks, not certified bounds on a continuum supremum. No measured fluid data or
full nonlinear water-wave simulation is added in this batch.

\textbf{Research connection.} Frequency targeting is motivated by Ben-Artzi,
Koch and Saut, \href{https://doi.org/10.1081/PDE-120025491}{\emph{Dispersion Estimates
for Third Order Equations in Two Dimensions}} (2003). Stationary-point-aware
integration is motivated by Huybrechs, Kuijlaars and Lejon,
\href{https://arxiv.org/abs/1806.06549}{\emph{A Numerical Method for Oscillatory
Integrals with Coalescing Saddle Points}} (2019), and Gibbs, Hewett and Huybrechs,
\href{https://arxiv.org/abs/2307.07261}{\emph{Numerical evaluation of oscillatory
integrals via automated steepest descent contour deformation}} (2024).
Our adaptive method is a simpler comparison calculation; it does not implement
their published algorithms or claim their error guarantees.

\textbf{Reproduce.} Run \texttt{research\_experiments.py} with arguments
\texttt{scaling}, \texttt{sweep}, \texttt{quadrature}, and \texttt{identity}; then
run \texttt{refine\_research\_flags.py} and \texttt{report/make\_research\_report.py}.
Case-level checkpoints, configurations and summaries are in
\texttt{results/research\_experiments/}. \par The annotated reading note is \texttt{research/literature\_review.md}.
